\section{Related work and contribution}
\label{sec:literature}

The empirical literature on minimum-wage price pass-through is largely built on
settings with substantial subnational policy variation, most commonly the United
States, where state-level minima diverge from the federal floor for extended
periods and across many states. That variation is what makes exposure-based
designs feasible. Portugal has been studied mainly through national time-series
and firm-level employment margins.

Our contribution is deliberately modest in method and specific in substance.

First, we assemble the statutory series from primary law rather than from
secondary compilations. This is not a formality. The published official history
omits the act effective in 2000 while the entries for 2001 state their increase
relative to it; three legally distinct minimum wages coexisted until 1991 and
two until 2004, so a single pre-2005 ``Portuguese minimum wage'' is a choice
that should be stated rather than assumed; and the two autonomous regions set
their wages by mechanisms that differ in a way that matters for identification.
Each of these is documented and each value is read from the act that set it.

Second, we provide the long-run accounting layer for Portugal on a consistent
basis from 1974, which requires reaching past the Eurostat series that begin in
the mid-1990s.

Third, we document a negative identification result. We are not aware of a
statement in the literature that the Azorean regional supplement, being
proportional, contributes no independent variation to a log-linear design. It is
a small point, and it removes one of the two regions that a researcher would
naturally expect to use.

We make no claim to be the first to study Portuguese minimum wages, and no claim
about causal pass-through.
